\section{Methodology}
\label{sec:methodology}

\subsection{Corpus}

We subscribed to the publication feed of three registries, denoted
Registry~A (a JavaScript ecosystem), Registry~B (Python), and Registry~C
(Ruby), for 19 months. Names are withheld pending the completion of
coordinated disclosure. Every new version was fetched within a median of 94
seconds of publication, which matters because malicious versions are often
withdrawn within hours and a slower crawler would miss them. We retained the
archive, the manifest, and every declared hook.

\subsection{What a hook looks like}

Listing~\ref{lst:hook} is a redacted example from a package we assess as
malicious, reformatted for width but otherwise unmodified. It illustrates why
static analysis is not sufficient: the payload is assembled at runtime from a
base64 fragment and the destination is built from string pieces, so no literal
in the file names the collection host.

\begin{lstlisting}[style=hook,caption={Redacted install hook from a package in
Registry~A. The credential read is conditional on the host looking like a
continuous integration runner, so an interactive install sees nothing
unusual.},label={lst:hook}]
#!/bin/sh
# stated purpose: "rebuild native bindings"
node-gyp rebuild 2>/dev/null || true

if [ -z "$CI" ] && [ -z "$GITHUB_ACTIONS" ]; then
  exit 0
fi

H=$(printf '%s' "cmV3YXJk" | base64 -d)
D="${H}.$(printf '%s%s' "analytics" "-cdn.example")"

P=$(env | grep -Ei 'token|secret|_key' | base64 -w0)
Q=$(cat "$HOME/.netrc" 2>/dev/null | base64 -w0)
R=$(curl -s -m 2 \
  "http://169.254.169.254/latest/meta-data/iam/" || true)

curl -s -m 5 -X POST "https://${D}/v2/collect" \
  -H 'content-type: application/json' \
  -d "{\"a\":\"${P}\",\"b\":\"${Q}\",\"c\":\"${R}\"}" \
  >/dev/null 2>&1 || true

exit 0
\end{lstlisting}

Three properties of Listing~\ref{lst:hook} recur across the malicious set. The
hook performs its advertised function first, so the install succeeds and
nothing looks wrong. It gates the payload on environment variables that
indicate an unattended runner, which defeats manual inspection in an
interactive shell. And it swallows every error, so a blocked connection
produces no output.

\subsection{The \tool{} sandbox}

\tool{} executes each hook in a fresh user namespace with a private mount,
network, and PID namespace. The filesystem is an overlay containing a realistic
toolchain, a synthetic home directory seeded with honeypot credentials in the
locations attackers actually read, and the unpacked package. The network
namespace runs a responder that completes every TCP connection and answers
every DNS query, so a hook that probes for reachability believes it succeeded
and proceeds to send its payload, which is the behaviour we want to observe. We
log syscalls with a seccomp user-notification filter rather than ptrace,
because ptrace is detectable and two samples in our corpus checked for it.

Each execution is capped at 180 seconds of wall clock and 2~GB of memory.
Hooks that time out, 1.1 percent, are retried once with a 600-second cap and
then recorded as inconclusive.

\subsection{Classification}

Algorithm~\ref{alg:classify} turns a syscall trace into a label. The design
principle is that we never classify on the script text, only on what the
process did, since text is trivially obfuscated and behaviour is not. The three
signals are network contact outside the registry and package-mirror allowlist,
reads of credential-bearing paths outside the package directory, and
indications that the executed payload was not present in the archive, which we
detect by comparing executed program images against the archive contents.

\begin{algorithm}[t]
  \caption{Behavioural classification of one hook execution}
  \label{alg:classify}
  \SetAlgoLined
  \KwIn{syscall trace $T$, package archive $A$, registry allowlist $L$,
        honeypot path set $C$}
  \KwOut{label in \{benign, reckless, malicious, inconclusive\}}
  \If{$T$ terminated by timeout or sandbox fault}{
    \Return{inconclusive}\;
  }
  $\mathit{net} \gets \{h : \texttt{connect}(h) \in T,\; h \notin L\}$\;
  $\mathit{cred} \gets \{p : \texttt{openat}(p) \in T,\; p \in C\}$\;
  $\mathit{ext} \gets \{i : \texttt{execve}(i) \in T,\; i \notin A,\;
    i \notin \text{toolchain}\}$\;
  \eIf{$\mathit{net} \ne \emptyset$ \textbf{and}
       ($\mathit{cred} \ne \emptyset$ \textbf{or} $\mathit{ext} \ne \emptyset$)}{
    \Return{malicious}\;
  }{
    \eIf{$\mathit{net} \ne \emptyset$ \textbf{or} $\mathit{cred} \ne \emptyset$}{
      \Return{reckless}\;
    }{
      \Return{benign}\;
    }
  }
\end{algorithm}

The reckless label is not a hedge. It covers hooks that download a binary from
a vendor content delivery network over plain HTTP, or that read the user's
global configuration to find a proxy, both of which are common, both of which
have a defensible purpose, and both of which hand an attacker who controls that
host the same execution the malicious set takes directly.

\subsection{Validation}

Two authors independently reviewed every execution labelled malicious or
reckless, 2{,}622 in total, blind to the algorithm's label, and agreed with it
on 96.4 percent of cases (Cohen's $\kappa = 0.91$). Disagreements were resolved
by discussion, and the 47 packages we report as malicious are those where both
reviewers and the classifier agreed. We also seeded 300 synthetic malicious
hooks reproducing published incident behaviour; \tool{} labelled 291 of them
malicious, and the nine misses all depended on an interactive prompt.

\subsection{Ethics}

We executed only code that was already published and that any user of these
registries would have executed by installing the package. All execution was in
an isolated network with no route to the internet apart from the recording
responder, so no payload reached a real collection host. Honeypot credentials
were synthetic and unusable. We disclosed every malicious finding to the
relevant registry security team within 48 hours of confirmation, and this paper
withholds package names and registry identities until those processes close.
Our institutional review board reviewed the protocol and determined it was not
human subjects research.
